% PDF 329--335；原书第 301--307 页。
\subsection{评估 RAG 系统能力}
除了用定量指标直接评估 RAG 系统的检索质量、答案忠实度和相关性，还需要考察几项能力，以了解系统中大语言模型生成器的总体表现、适用范围和可靠程度（\refcite{ch07-ref-002}{Chen et al., 2023a}）。这类评估有助于弄清检索上下文和提示扩充怎样影响模型生成。例如，要了解生成器在多大程度上采用非参数知识，即提示中的上下文信息，替代参数知识，即预训练时存入模型参数的信息。如果生成器用检索到的非参数知识覆盖参数知识，而检索结果有误，开发人员就可能需要修改提示，或提高检索精确率，使生成的回答尽量符合事实。

下面介绍常用的四项 RAG 能力评估。图例取自 \refcite{ch07-ref-002}{Chen 等人（2023a）}，图~\ref{fig:ch07-capabilities} 汇总各项能力。

\subsubsection{噪声鲁棒性}
噪声鲁棒性（noise robustness）衡量大语言模型生成器能否在含有噪声的检索文档中，只采用有用信息，如图~\ref{fig:ch07-noise} 所示。它考察生成器在上下文混有无关信息时，能否给出用户查询的正确答案。

\sourcefigure{fig07_08.pdf}{RAG 回答中的噪声鲁棒性示例。虽然检索上下文含有噪声，例如关于诺贝尔物理学奖而非文学奖的信息，生成器仍能给出用户查询的正确答案。}{fig:ch07-noise}

\begin{practicaltip}
要评估 RAG 大语言模型处理含噪上下文的能力，需要知道对于一组生成的问答对，哪些上下文相关，哪些不相关。通常将相关文档与按适当比例随机抽取的不相关上下文配对。再将这些“检索”上下文交给生成器，并将它的答案与“问题—支持信息—答案”三元组数据集中的原答案比较。可用精确匹配（exact match，EM）等准确性指标量化模型在这方面的能力（\refcite{ch07-ref-002}{Chen et al., 2023a}）。
\end{practicaltip}

\subsubsection{无答案时拒答}
在无答案时拒答（negative rejection）这一能力中，如果没有任何检索上下文包含回答用户查询所需的相关信息，RAG 应用就应拒绝作答。如图~\ref{fig:ch07-rejection} 所示，面对“谁获得了 2022 年诺贝尔文学奖？”这一问题，所有上下文都缺少回答所需的事实，取回的内容仅与诺贝尔物理学奖有关。

评估这种能力，有助于开发人员改进应用在知识来源不足以支持忠实或事实正确的回答时的行为。例如，可以针对这种情况设置更明确的系统指令。%
\footnote{译注：原书图 7.9 的图注写作 2021 年，图内的用户查询和正文都写作 2022 年。}

\sourcefigure{fig07_09.pdf}{没有任何检索上下文包含回答所需信息时，如果生成器不给出用户查询的答案，就表现出无答案时拒答的能力。本例的上下文没有包含 2022 年诺贝尔文学奖得主的信息，因此生成器表示无法回答。}{fig:ch07-rejection}

\begin{practicaltip}
评估无答案时拒答的能力，同样先生成一组问答对及支持答案的信息。然后只从不相关文档抽取上下文，即抽取那些并非随问答对生成的文档。根据模型正确给出特定拒答内容的次数，可衡量 RAG 大语言模型的拒答能力。拒答内容通常由给模型的指令指定。计算拒答率，可了解查询 RAG 系统时正确和错误拒答各占多少。
\end{practicaltip}

\subsubsection{信息整合}
RAG 中的信息整合（information integration）指大语言模型生成器能从多份上下文文档中综合信息，为复杂问题给出正确答案。如果用户查询包含多个子问题，或涉及分布在各处的多项信息，RAG 系统需要将各项信息对应的相关上下文整合进回答。如图~\ref{fig:ch07-integration} 所示，用户询问 2021 年和 2022 年诺贝尔物理学奖的得主。回答所需的外部信息也分布在两份上下文文档中。RAG 回答正确地综合了这些上下文，给出了所需的答案。

\sourcefigure{fig07_10.pdf}{RAG 生成器能利用多份上下文文档回答复杂问题时，表现出信息整合能力。本例中，用户同时询问 2021 年和 2022 年诺贝尔物理学奖的得主，生成器综合相关上下文，给出正确答案。}{fig:ch07-integration}

\begin{practicaltip}
实践中，评估信息整合也依赖于生成“问题—支持信息—答案”三元组。但在构造测试数据时，要多做一步，为问题增加应回答的方面。例如，可以合并两个问题及其各自的答案和支持信息，使回答这个较复杂的问题所需的信息分散在多份上下文文档中。也可用精确匹配等准确性指标判断信息是否成功整合，直接比较 RAG 生成的回答与原先为问题生成的答案（\refcite{ch07-ref-002}{Chen et al., 2023a}）。
\end{practicaltip}

\subsubsection{反事实鲁棒性}
RAG 应用常用的外部知识库也可能含有事实错误。因此，需要评估 RAG 生成器能否识别检索上下文中的错误，这称为反事实鲁棒性（counterfactual robustness）。识别检索上下文中的错误，完全依赖生成器的参数知识。如果所选大语言模型没有相应领域的知识，或掌握得不充分，这项能力就难以评估。可以对大语言模型做领域适配或微调，但费用较高，也会抵消 RAG 的一些优点。不过，开源领域已有许多经过领域微调的大语言模型，因此生成适用的领域测试数据也变得更可行。

\sourcefigure{fig07_11.pdf}{反事实鲁棒性指生成器能检测所提供上下文的事实错误，并在回答中指出错误。这项能力依靠大语言模型生成器的参数知识。如果模型缺少相应领域知识，或应用依赖模型知识截止日期之后的信息，这项能力可能难以评估。}{fig:ch07-counterfactual}

\begin{practicaltip}
检验这项能力时，先提示大语言模型生成器只根据其参数知识生成问题和答案。也就是说，让模型提出它已经知道答案的问题，无需依赖外部知识库。为构造反事实支持信息，先核实生成的答案正确，再从外部知识中检索支持信息。然后手工修改支持上下文，用错误的事实代替正确的支持事实。用这个数据集考察两项指标，即错误检测率和错误纠正率。

错误检测率取决于支持上下文含有事实错误时，生成器是否返回指定内容。它衡量大语言模型能在多大程度上根据参数知识判断检索上下文的事实性。类似地，错误纠正率衡量支持信息包含错误时，生成器仍能给出正确答案的频率。
\end{practicaltip}

\begin{figure}[!htbp]
\centering
% 重绘自原书第 306 页的图 7.12。
\begin{tikzpicture}[
  font=\footnotesize\sffamily,
  topbox/.style={fill=yellow!26,minimum width=24mm,minimum height=8mm,
    inner sep=2pt,align=center},
  band/.style={draw=black!78,minimum width=109mm,minimum height=8.7mm,
    align=center,font=\sffamily\bfseries},
  line/.style={draw=black!71,line width=.56pt},
  arrow/.style={line,-{Stealth[length=1.5mm]}}]
\node[band,fill=violet!17] (response) at (0,0) {生成回答};
\node[band,fill=green!18] (documents) at (0,-3.67) {检索文档};
\foreach \x/\item in {-3.99/噪声鲁棒性,-1.33/无答案时拒答,1.33/信息整合,3.99/反事实鲁棒性}{
  \node[topbox] (metric) at (\x,2.32) {\item};
  \draw[line] (metric.south) -- (response.north -| metric);
}
\node[fill=white,inner sep=2pt,font=\scriptsize\sffamily] at (0,1.15)
  {衡量生成回答的下列能力};
\foreach \x/\y/\item in {
   -3.99/-1.36/忽略其中的\\无关信息,
   -1.33/-2.44/承认其中缺少\\回答所需的信息,
   1.33/-1.36/整合来自\\多份文档的信息,
   3.99/-2.44/识别并拒绝\\其中的错误信息}{
  \coordinate (lane) at (\x,0);
  \draw[arrow] (response.south -| lane) -- (documents.north -| lane);
  \node[fill=white,inner sep=1.8pt,align=center,font=\scriptsize\sffamily] at (\x,\y) {\item};
}
\end{tikzpicture}

\caption{RAG 系统各项能力指标所衡量的性质。四项指标各自衡量生成器能否理解检索文档所具有的正面或负面性质，并作出合适的回答。}
\label{fig:ch07-capabilities}
\end{figure}

\subsection{RAG 评估小结}
RAG 应用有多个边界不容易划清的功能组件，因此评估较为复杂。开发人员需要考虑向量数据库的索引是否合适、外部知识文档如何分块、检索结果的相关性与排序，以及最后的生成阶段。已有许多研究与工具简化了这个过程。从 TruLens 团队提出的 RAG 三项指标框架，到 LlamaIndex 等便于实现的框架，这些工具减少了需要用户自行处理的细节，使其能够快速制作原型，并开发适合生产使用的应用。

第~8~章将讨论在生产环境中使用大语言模型时涉及的运维概念、框架、工具与挑战，其中不少也适用于 RAG 应用开发。深入讨论这些问题之前，先给出一个 RAG 开发和评估的教程。
